\chapter{Detailed Oracle Partitioning Traces}
\label{app:oracle-partition-details}

This appendix provides the edge-level evidence and algorithm traces underlying the
oracle partitions summarized in Chapter~\ref{sec:results}. The New German Credit
trace uses Cram\'er's $V$ as its edge weight, whereas the natural-language RQ5
example uses the empirical probability that two concepts were placed in the same
group across repeated auxiliary-LLM judgments.

\section{New German Credit}
\label{app:credit-partition-trace}

Figures~\ref{fig:credit-grouping-steps-a} and~\ref{fig:credit-grouping-steps-b}
show the complete execution of the constrained greedy partitioning algorithm at
$n_{\max}=3$. The initialization contains twelve singleton components and all
retained candidate edges. Each subsequent panel adds one accepted edge in descending
order of association strength, provided that merging its components does not exceed
the component-size constraint. Green edges were accepted in an earlier step, the
orange edge is the edge accepted in the displayed step, and gray edges remain
unselected. Thus, the seven accepted edges produce the five final components shown
in Figure~\ref{fig:german-credit-grouping-final}.

% Reusable panel renderer for the New German Credit greedy-merge trace.
% #1: number of accepted merge steps already processed (0--7)
% #2: panel title
\definecolor{creditTraceAccepted}{HTML}{17823B}
\newcommand{\CreditGroupingStage}[2]{%
  \begin{tikzpicture}[
    x=0.47cm,
    y=0.47cm,
    credit stage node/.style={
      draw=black!60,
      fill=black!7,
      rounded corners=1.5pt,
      line width=0.45pt,
      minimum height=5.5mm,
      text width=15mm,
      align=center,
      inner sep=1.5pt,
      font=\sffamily\tiny\bfseries
    },
    pending stage edge/.style={draw=black!22,line width=0.35pt},
    accepted stage edge/.style={draw=creditTraceAccepted,line width=1.35pt},
    current stage edge/.style={draw=orange!85!black,line width=1.8pt}
  ]
    \node[font=\small\bfseries,anchor=south] at (0.6,4.75) {#2};

    \newcommand{\CreditNode}[5]{%
      \node[credit stage node] (##1) at (##3,##4) {##2};%
    }
    \newcommand{\CreditEdge}[5]{}%
    % Data for the New German Credit association graph.
%
% Source: FELICS data/new_german_credit/data.csv, processed as in
% src/intervention_generation/tabular_intervention_generator.py with
% p < 0.01, Cramer's V >= 0.5, and n_max = 3. Values are rounded only for
% display. The final argument of \CreditEdge is the greedy merge step; zero
% marks a retained association that was not selected for the spanning forest.
%
% Keeping the records separate from the renderer allows the same graph to be
% reused with another \CreditGroupingStep or alternative edge styles.

% \CreditNode{id}{label}{x}{y}{final group}
\CreditNode{income}{Household Income}{-0.8}{4.0}{One}
\CreditNode{liquid}{Liquid Assets}{2.7}{3.4}{One}
\CreditNode{schufa}{SCHUFA Score (\%)}{1.1}{-2.4}{One}

\CreditNode{amount}{Credit Amount}{3.9}{0.9}{Two}
\CreditNode{purpose}{Credit Purpose}{1.5}{-0.2}{Two}
\CreditNode{term}{Credit Term (Months)}{3.8}{-3.8}{Two}

\CreditNode{realestate}{Real Estate Value}{-0.9}{0.8}{Three}
\CreditNode{ownhouse}{Lives in Own House}{3.0}{-1.4}{Three}
\CreditNode{age}{Applicant Age}{-2.1}{-3.8}{Three}

\CreditNode{job}{Job Category}{-4.0}{2.4}{Four}
\CreditNode{debt}{Monthly Debt Burden}{-5.0}{-1.3}{Four}

\CreditNode{householdsize}{Household Size}{6.2}{1.2}{Five}

% \CreditEdge{first node}{second node}{Cramer's V}{merge step}{label position}
\CreditEdge{realestate}{ownhouse}{0.90}{1}{0.55}
\CreditEdge{income}{liquid}{0.71}{2}{0.54}
\CreditEdge{amount}{term}{0.68}{3}{0.55}
\CreditEdge{amount}{purpose}{0.65}{4}{0.48}
\CreditEdge{debt}{job}{0.61}{5}{0.48}
\CreditEdge{schufa}{liquid}{0.60}{6}{0.48}
\CreditEdge{age}{ownhouse}{0.59}{7}{0.55}

\CreditEdge{term}{purpose}{0.64}{0}{0.46}
\CreditEdge{income}{amount}{0.61}{0}{0.55}
\CreditEdge{income}{job}{0.59}{0}{0.46}
\CreditEdge{purpose}{job}{0.58}{0}{0.52}
\CreditEdge{job}{ownhouse}{0.57}{0}{0.52}
\CreditEdge{purpose}{realestate}{0.57}{0}{0.47}
\CreditEdge{liquid}{realestate}{0.56}{0}{0.48}
\CreditEdge{income}{purpose}{0.56}{0}{0.47}
\CreditEdge{age}{purpose}{0.56}{0}{0.52}
\CreditEdge{amount}{liquid}{0.55}{0}{0.43}
\CreditEdge{age}{realestate}{0.55}{0}{0.52}
\CreditEdge{purpose}{liquid}{0.55}{0}{0.45}
\CreditEdge{job}{realestate}{0.54}{0}{0.54}
\CreditEdge{job}{liquid}{0.53}{0}{0.48}
\CreditEdge{purpose}{ownhouse}{0.52}{0}{0.53}
\CreditEdge{income}{realestate}{0.52}{0}{0.52}
\CreditEdge{income}{schufa}{0.50}{0}{0.58}


    \begin{scope}[on background layer]
      \renewcommand{\CreditNode}[5]{}%
      \renewcommand{\CreditEdge}[5]{%
        \ifnum##4=0
          \draw[pending stage edge] (##1) -- (##2);%
        \else
          \ifnum##4>#1
            \draw[pending stage edge] (##1) -- (##2);%
          \else
            \ifnum##4=#1
              \draw[current stage edge] (##1) -- (##2);%
            \else
              \draw[accepted stage edge] (##1) -- (##2);%
            \fi
          \fi
        \fi
      }
      % Data for the New German Credit association graph.
%
% Source: FELICS data/new_german_credit/data.csv, processed as in
% src/intervention_generation/tabular_intervention_generator.py with
% p < 0.01, Cramer's V >= 0.5, and n_max = 3. Values are rounded only for
% display. The final argument of \CreditEdge is the greedy merge step; zero
% marks a retained association that was not selected for the spanning forest.
%
% Keeping the records separate from the renderer allows the same graph to be
% reused with another \CreditGroupingStep or alternative edge styles.

% \CreditNode{id}{label}{x}{y}{final group}
\CreditNode{income}{Household Income}{-0.8}{4.0}{One}
\CreditNode{liquid}{Liquid Assets}{2.7}{3.4}{One}
\CreditNode{schufa}{SCHUFA Score (\%)}{1.1}{-2.4}{One}

\CreditNode{amount}{Credit Amount}{3.9}{0.9}{Two}
\CreditNode{purpose}{Credit Purpose}{1.5}{-0.2}{Two}
\CreditNode{term}{Credit Term (Months)}{3.8}{-3.8}{Two}

\CreditNode{realestate}{Real Estate Value}{-0.9}{0.8}{Three}
\CreditNode{ownhouse}{Lives in Own House}{3.0}{-1.4}{Three}
\CreditNode{age}{Applicant Age}{-2.1}{-3.8}{Three}

\CreditNode{job}{Job Category}{-4.0}{2.4}{Four}
\CreditNode{debt}{Monthly Debt Burden}{-5.0}{-1.3}{Four}

\CreditNode{householdsize}{Household Size}{6.2}{1.2}{Five}

% \CreditEdge{first node}{second node}{Cramer's V}{merge step}{label position}
\CreditEdge{realestate}{ownhouse}{0.90}{1}{0.55}
\CreditEdge{income}{liquid}{0.71}{2}{0.54}
\CreditEdge{amount}{term}{0.68}{3}{0.55}
\CreditEdge{amount}{purpose}{0.65}{4}{0.48}
\CreditEdge{debt}{job}{0.61}{5}{0.48}
\CreditEdge{schufa}{liquid}{0.60}{6}{0.48}
\CreditEdge{age}{ownhouse}{0.59}{7}{0.55}

\CreditEdge{term}{purpose}{0.64}{0}{0.46}
\CreditEdge{income}{amount}{0.61}{0}{0.55}
\CreditEdge{income}{job}{0.59}{0}{0.46}
\CreditEdge{purpose}{job}{0.58}{0}{0.52}
\CreditEdge{job}{ownhouse}{0.57}{0}{0.52}
\CreditEdge{purpose}{realestate}{0.57}{0}{0.47}
\CreditEdge{liquid}{realestate}{0.56}{0}{0.48}
\CreditEdge{income}{purpose}{0.56}{0}{0.47}
\CreditEdge{age}{purpose}{0.56}{0}{0.52}
\CreditEdge{amount}{liquid}{0.55}{0}{0.43}
\CreditEdge{age}{realestate}{0.55}{0}{0.52}
\CreditEdge{purpose}{liquid}{0.55}{0}{0.45}
\CreditEdge{job}{realestate}{0.54}{0}{0.54}
\CreditEdge{job}{liquid}{0.53}{0}{0.48}
\CreditEdge{purpose}{ownhouse}{0.52}{0}{0.53}
\CreditEdge{income}{realestate}{0.52}{0}{0.52}
\CreditEdge{income}{schufa}{0.50}{0}{0.58}

    \end{scope}
  \end{tikzpicture}%
}


\begin{figure}[p]
  \centering
  \begin{minipage}[t]{0.49\textwidth}\centering
    \CreditGroupingStage{0}{Initialization}
  \end{minipage}\hfill
  \begin{minipage}[t]{0.49\textwidth}\centering
    \CreditGroupingStage{1}{Step 1: Real Estate--Own House}
  \end{minipage}

  \vspace{1em}
  \begin{minipage}[t]{0.49\textwidth}\centering
    \CreditGroupingStage{2}{Step 2: Income--Liquid Assets}
  \end{minipage}\hfill
  \begin{minipage}[t]{0.49\textwidth}\centering
    \CreditGroupingStage{3}{Step 3: Amount--Term}
  \end{minipage}
  \caption{Initialization and the first three accepted merges of the New German Credit partitioning algorithm.}
  \label{fig:credit-grouping-steps-a}
\end{figure}

\begin{figure}[p]
  \centering
  \begin{minipage}[t]{0.49\textwidth}\centering
    \CreditGroupingStage{4}{Step 4: Amount--Purpose}
  \end{minipage}\hfill
  \begin{minipage}[t]{0.49\textwidth}\centering
    \CreditGroupingStage{5}{Step 5: Debt Burden--Job}
  \end{minipage}

  \vspace{1em}
  \begin{minipage}[t]{0.49\textwidth}\centering
    \CreditGroupingStage{6}{Step 6: SCHUFA--Liquid Assets}
  \end{minipage}\hfill
  \begin{minipage}[t]{0.49\textwidth}\centering
    \CreditGroupingStage{7}{Step 7: Age--Own House}
  \end{minipage}
  \caption{The final four accepted merges of the New German Credit partitioning algorithm. The last panel is the completed maximum spanning forest.}
  \label{fig:credit-grouping-steps-b}
\end{figure}
\FloatBarrier

\begingroup
\small
\newcommand{\CreditNode}[5]{%
  \expandafter\def\csname CreditAppendixNodeLabel@#1\endcsname{#2}%
}
\newcommand{\CreditEdge}[5]{}
% Data for the New German Credit association graph.
%
% Source: FELICS data/new_german_credit/data.csv, processed as in
% src/intervention_generation/tabular_intervention_generator.py with
% p < 0.01, Cramer's V >= 0.5, and n_max = 3. Values are rounded only for
% display. The final argument of \CreditEdge is the greedy merge step; zero
% marks a retained association that was not selected for the spanning forest.
%
% Keeping the records separate from the renderer allows the same graph to be
% reused with another \CreditGroupingStep or alternative edge styles.

% \CreditNode{id}{label}{x}{y}{final group}
\CreditNode{income}{Household Income}{-0.8}{4.0}{One}
\CreditNode{liquid}{Liquid Assets}{2.7}{3.4}{One}
\CreditNode{schufa}{SCHUFA Score (\%)}{1.1}{-2.4}{One}

\CreditNode{amount}{Credit Amount}{3.9}{0.9}{Two}
\CreditNode{purpose}{Credit Purpose}{1.5}{-0.2}{Two}
\CreditNode{term}{Credit Term (Months)}{3.8}{-3.8}{Two}

\CreditNode{realestate}{Real Estate Value}{-0.9}{0.8}{Three}
\CreditNode{ownhouse}{Lives in Own House}{3.0}{-1.4}{Three}
\CreditNode{age}{Applicant Age}{-2.1}{-3.8}{Three}

\CreditNode{job}{Job Category}{-4.0}{2.4}{Four}
\CreditNode{debt}{Monthly Debt Burden}{-5.0}{-1.3}{Four}

\CreditNode{householdsize}{Household Size}{6.2}{1.2}{Five}

% \CreditEdge{first node}{second node}{Cramer's V}{merge step}{label position}
\CreditEdge{realestate}{ownhouse}{0.90}{1}{0.55}
\CreditEdge{income}{liquid}{0.71}{2}{0.54}
\CreditEdge{amount}{term}{0.68}{3}{0.55}
\CreditEdge{amount}{purpose}{0.65}{4}{0.48}
\CreditEdge{debt}{job}{0.61}{5}{0.48}
\CreditEdge{schufa}{liquid}{0.60}{6}{0.48}
\CreditEdge{age}{ownhouse}{0.59}{7}{0.55}

\CreditEdge{term}{purpose}{0.64}{0}{0.46}
\CreditEdge{income}{amount}{0.61}{0}{0.55}
\CreditEdge{income}{job}{0.59}{0}{0.46}
\CreditEdge{purpose}{job}{0.58}{0}{0.52}
\CreditEdge{job}{ownhouse}{0.57}{0}{0.52}
\CreditEdge{purpose}{realestate}{0.57}{0}{0.47}
\CreditEdge{liquid}{realestate}{0.56}{0}{0.48}
\CreditEdge{income}{purpose}{0.56}{0}{0.47}
\CreditEdge{age}{purpose}{0.56}{0}{0.52}
\CreditEdge{amount}{liquid}{0.55}{0}{0.43}
\CreditEdge{age}{realestate}{0.55}{0}{0.52}
\CreditEdge{purpose}{liquid}{0.55}{0}{0.45}
\CreditEdge{job}{realestate}{0.54}{0}{0.54}
\CreditEdge{job}{liquid}{0.53}{0}{0.48}
\CreditEdge{purpose}{ownhouse}{0.52}{0}{0.53}
\CreditEdge{income}{realestate}{0.52}{0}{0.52}
\CreditEdge{income}{schufa}{0.50}{0}{0.58}

\renewcommand{\CreditNode}[5]{}
\renewcommand{\CreditEdge}[5]{%
  \csname CreditAppendixNodeLabel@#1\endcsname &
  \csname CreditAppendixNodeLabel@#2\endcsname & #3\tabularnewline
}
\begin{longtable}{@{}p{0.40\textwidth}p{0.40\textwidth}>{\raggedleft\arraybackslash}p{0.10\textwidth}@{}}
  \caption{Retained pairwise associations used for the New German Credit partition.}
  \label{tab:credit-associations}\\
  \hline
  \textbf{Node 1} & \textbf{Node 2} & \textbf{$V$}\\
  \hline
  \endfirsthead
  \hline
  \textbf{Node 1} & \textbf{Node 2} & \textbf{$V$}\\
  \hline
  \endhead
  % Data for the New German Credit association graph.
%
% Source: FELICS data/new_german_credit/data.csv, processed as in
% src/intervention_generation/tabular_intervention_generator.py with
% p < 0.01, Cramer's V >= 0.5, and n_max = 3. Values are rounded only for
% display. The final argument of \CreditEdge is the greedy merge step; zero
% marks a retained association that was not selected for the spanning forest.
%
% Keeping the records separate from the renderer allows the same graph to be
% reused with another \CreditGroupingStep or alternative edge styles.

% \CreditNode{id}{label}{x}{y}{final group}
\CreditNode{income}{Household Income}{-0.8}{4.0}{One}
\CreditNode{liquid}{Liquid Assets}{2.7}{3.4}{One}
\CreditNode{schufa}{SCHUFA Score (\%)}{1.1}{-2.4}{One}

\CreditNode{amount}{Credit Amount}{3.9}{0.9}{Two}
\CreditNode{purpose}{Credit Purpose}{1.5}{-0.2}{Two}
\CreditNode{term}{Credit Term (Months)}{3.8}{-3.8}{Two}

\CreditNode{realestate}{Real Estate Value}{-0.9}{0.8}{Three}
\CreditNode{ownhouse}{Lives in Own House}{3.0}{-1.4}{Three}
\CreditNode{age}{Applicant Age}{-2.1}{-3.8}{Three}

\CreditNode{job}{Job Category}{-4.0}{2.4}{Four}
\CreditNode{debt}{Monthly Debt Burden}{-5.0}{-1.3}{Four}

\CreditNode{householdsize}{Household Size}{6.2}{1.2}{Five}

% \CreditEdge{first node}{second node}{Cramer's V}{merge step}{label position}
\CreditEdge{realestate}{ownhouse}{0.90}{1}{0.55}
\CreditEdge{income}{liquid}{0.71}{2}{0.54}
\CreditEdge{amount}{term}{0.68}{3}{0.55}
\CreditEdge{amount}{purpose}{0.65}{4}{0.48}
\CreditEdge{debt}{job}{0.61}{5}{0.48}
\CreditEdge{schufa}{liquid}{0.60}{6}{0.48}
\CreditEdge{age}{ownhouse}{0.59}{7}{0.55}

\CreditEdge{term}{purpose}{0.64}{0}{0.46}
\CreditEdge{income}{amount}{0.61}{0}{0.55}
\CreditEdge{income}{job}{0.59}{0}{0.46}
\CreditEdge{purpose}{job}{0.58}{0}{0.52}
\CreditEdge{job}{ownhouse}{0.57}{0}{0.52}
\CreditEdge{purpose}{realestate}{0.57}{0}{0.47}
\CreditEdge{liquid}{realestate}{0.56}{0}{0.48}
\CreditEdge{income}{purpose}{0.56}{0}{0.47}
\CreditEdge{age}{purpose}{0.56}{0}{0.52}
\CreditEdge{amount}{liquid}{0.55}{0}{0.43}
\CreditEdge{age}{realestate}{0.55}{0}{0.52}
\CreditEdge{purpose}{liquid}{0.55}{0}{0.45}
\CreditEdge{job}{realestate}{0.54}{0}{0.54}
\CreditEdge{job}{liquid}{0.53}{0}{0.48}
\CreditEdge{purpose}{ownhouse}{0.52}{0}{0.53}
\CreditEdge{income}{realestate}{0.52}{0}{0.52}
\CreditEdge{income}{schufa}{0.50}{0}{0.58}

\end{longtable}
\endgroup
\clearpage

\section{RQ5 Motivating MedQA Example}
\label{app:rq5-oracle-partition}

For the controlled version of MedQA question 521, the natural-language oracle was
sampled 50 times. The weight of an edge is the proportion of these valid judgments
in which its two concepts occurred in the same candidate interaction group. Edges
with a co-grouping probability of at least $0.50$ were retained and processed by the
same constrained greedy merging procedure used for the tabular oracle. Figure~\ref{fig:rq5-oracle-grouping}
shows the resulting partition. The edge between refusal of treatment and
self-perception was retained but did not need to be accepted because both concepts
already belonged to the component connected through eating disorder.

\begin{figure}[H]
  \centering
  \includestandalone[mode=buildnew,width=0.78\textwidth,height=0.36\textheight,keepaspectratio]{graphs/rq5_oracle_grouping}
  \caption{Final consensus partition for the RQ5 motivating MedQA example. Edge labels are empirical co-grouping probabilities across 50 auxiliary-LLM judgments.}
  \label{fig:rq5-oracle-grouping}
\end{figure}

\begingroup
\small
\newcommand{\RQFiveNode}[5]{%
  \expandafter\def\csname RQFiveAppendixNodeLabel@#1\endcsname{#2}%
}
\newcommand{\RQFiveEdge}[4]{}
% Final consensus graph for controlled MedQA question 521 (RQ5).
%
% Source: experiment/intervention_generation/medqa/rq5_felics/controlled/
% medqa_521_matton_table9/example_521/correlations.json. The source records
% 50 valid auxiliary-LLM judgments and a retention threshold of 0.50.
%
% \RQFiveNode{id}{display label}{x}{y}{final group}
\RQFiveNode{age}{Age}{-4.1}{2.4}{Three}
\RQFiveNode{gender}{Gender}{-4.1}{0.2}{Four}
\RQFiveNode{mental}{Mental Status}{-2.7}{-2.2}{Five}
\RQFiveNode{reason}{Reason for Visit}{0.2}{-2.6}{Six}

\RQFiveNode{eating}{Eating Disorder}{-0.7}{2.5}{One}
\RQFiveNode{self}{Self-perception of Weight}{1.8}{2.3}{One}
\RQFiveNode{refusal}{Refusal of Treatment}{0.5}{0.3}{One}

\RQFiveNode{vitals}{Vital Signs}{3.8}{0.8}{Two}
\RQFiveNode{ems}{EMS Treatment}{4.4}{-1.7}{Two}

% \RQFiveEdge{first node}{second node}{co-grouping probability}{merge step}
\RQFiveEdge{eating}{self}{1.00}{1}
\RQFiveEdge{eating}{refusal}{0.86}{2}
\RQFiveEdge{refusal}{self}{0.86}{0}
\RQFiveEdge{vitals}{ems}{0.66}{3}

\renewcommand{\RQFiveNode}[5]{}
\renewcommand{\RQFiveEdge}[4]{%
  \csname RQFiveAppendixNodeLabel@#1\endcsname &
  \csname RQFiveAppendixNodeLabel@#2\endcsname & #3\tabularnewline
}
\begin{table}[H]
  \centering
  \caption{Retained co-grouping frequencies for the RQ5 motivating MedQA example.}
  \label{tab:rq5-cogrouping}
  \begin{tabular}{@{}p{0.38\textwidth}p{0.38\textwidth}>{\raggedleft\arraybackslash}p{0.15\textwidth}@{}}
    \hline
    \textbf{Node 1} & \textbf{Node 2} & \textbf{Co-grouping probability}\\
    \hline
    % Final consensus graph for controlled MedQA question 521 (RQ5).
%
% Source: experiment/intervention_generation/medqa/rq5_felics/controlled/
% medqa_521_matton_table9/example_521/correlations.json. The source records
% 50 valid auxiliary-LLM judgments and a retention threshold of 0.50.
%
% \RQFiveNode{id}{display label}{x}{y}{final group}
\RQFiveNode{age}{Age}{-4.1}{2.4}{Three}
\RQFiveNode{gender}{Gender}{-4.1}{0.2}{Four}
\RQFiveNode{mental}{Mental Status}{-2.7}{-2.2}{Five}
\RQFiveNode{reason}{Reason for Visit}{0.2}{-2.6}{Six}

\RQFiveNode{eating}{Eating Disorder}{-0.7}{2.5}{One}
\RQFiveNode{self}{Self-perception of Weight}{1.8}{2.3}{One}
\RQFiveNode{refusal}{Refusal of Treatment}{0.5}{0.3}{One}

\RQFiveNode{vitals}{Vital Signs}{3.8}{0.8}{Two}
\RQFiveNode{ems}{EMS Treatment}{4.4}{-1.7}{Two}

% \RQFiveEdge{first node}{second node}{co-grouping probability}{merge step}
\RQFiveEdge{eating}{self}{1.00}{1}
\RQFiveEdge{eating}{refusal}{0.86}{2}
\RQFiveEdge{refusal}{self}{0.86}{0}
\RQFiveEdge{vitals}{ems}{0.66}{3}

  \end{tabular}
\end{table}
\endgroup
